\section{Traversability, Motion Planning, and Control}
\label{sec:planning}

Motion planning and control connect the estimated environment to executable behavior. Under degradation, feasibility depends on terrain support, obstacle motion, state uncertainty, actuator limits, and the age of control information. The methods below address these dependencies through traversability costs, information-aware exploration, learned prediction and policies, safety constraints, and predictive control, with representative implementations listed in Table~\ref{S-tab:planning} of the supplementary material.

\subsection{Terrain Support and Traversability Costs}

Traversability is a relation between terrain and robot capability, which the reviewed methods express through quantities with different physical meanings. CADNet predicts heading-dependent accessible range from LiDAR history for a TEB planner on a SCOUT robot~\cite{P_M29}, whereas GA-Nav projects attention-based terrain segmentation into a cost map combined with elevation~\cite{P_M40}. A risk-aware learner labels traversed footprints and adds weak obstacle labels before predicting unvisited cells, although its predictions are evaluated only on manually driven recordings~\cite{P_M23}. PE-RRT* infers hidden support planes in vegetation with multivariate Gaussian processes and tracks the route with nonlinear MPC~\cite{P_M26}, while a multilevel planner checks settling, support polygon, and body collision for a tracked robot whose wheeled variant is simulated~\cite{P_M12}. An anisotropic explorer penalizes motion by slope and travel direction with a quasi-static Euler Cost that omits slip and suspension transients~\cite{P_N08}. A footprint label, a support plane, a stability margin, and a cross-slope penalty are not interchangeable estimates of traversal risk. Each is also computed from an estimated pose, and upstream pose error can therefore alter the planning input even when the terrain model is unchanged.

Recent planners condition terrain risk on approach direction and uncertainty and embed it in optimization. SPARTA predicts angle-conditioned risk distributions as reusable Fourier functions from a fully observed prior point cloud~\cite{P_S45}. Its tire-deformation consequences are evaluated in simulation, with a physical demonstration that routes around a virtual risky branch. A visual-foundation-model approach learns cost, speed, and novelty maps from expert driving for MPPI~\cite{P_S42}. In its field trials, operator corrections expose residual failures that offline trajectory agreement conceals. A terrain-adaptive $\epsilon$-constraint MPC minimizes path length while bounding a pose-stability cost adjusted online from terrain descriptors and corrects car-like dynamics with a sparse GP~\cite{P_R13}. Its success falls from 94\% in simulation to 80\% outdoors on a Scout Mini, a drop that the authors attribute to wheel slip and sensor noise, while the planner assumes an offline terrain mesh. TRAIL learns an implicit field of terrain attributes from RGB and LiDAR and optimizes path geometry and speed along its gradients, with field interventions caused by wet-ground traction loss and failures attributed to occlusion and limited field of view~\cite{P_R15}. These systems already make risk-sensitive terrain decisions, whereas map completeness, uncertainty calibration, traction, and physical consequences remain separate evaluation questions.

\subsection{Exploration with Information and Geometry}

Exploration chooses where to improve the map while preserving the ability to localize. FIT-SLAM~2 attaches frontiers to SLAM keyframes, adapts a traversability roadmap after pose changes, and rejects edges with low expected landmark information~\cite{P_E13}. With a common controller and simulator pose, completion takes 880\,s versus 1118\,s for FAEL in Cave but 1809\,s versus 1706\,s in Corridor. Fisher-information evaluation also raises mean planning time from 20.8 to 63.0\,ms. Confidence-rich exploration combines particle-trajectory entropy with occupancy confidence and executes the loop on hardware~\cite{P_E14}, whereas an uncertainty-frontier planner scores coverage with signed relative entropy under known orientation~\cite{P_E12}. Other explorers change the representation of unexplored space, including safety levels for unknown regions in LRAE~\cite{P_M14}, octree category distributions for semantic exploration~\cite{P_M31}, and overlapping free polytopes in FRTree, whose body-containment corridor does not enforce wheeled curvature or rolling-contact limits~\cite{P_M05}. A SLAM-aware planner inserts informative loop closures over a supplied topo-metric graph~\cite{P_N05}. Potential-field MPC integrates frontier and unknown-space costs on a six-wheel skid-steer robot, although without systematic baselines~\cite{P_S02}. The FIT-SLAM~2 results show that more observable routes can cost travel and computation. The benefit of estimation-aware planning further depends on the availability of priors and on useful loop opportunities.

\subsection{Learning, Prediction, and Interaction}

Learning enters planning at different points. NAMR-RRT learns where a risk-aware tree samples~\cite{P_M08}, prescient navigation passes multimodal pedestrian hypotheses to nonlinear MPC that may exhaust its time budget~\cite{P_M30}, and a learned relaxation of optimal reciprocal collision avoidance (ORCA) predicts escape direction and avoidance responsibility~\cite{P_M19}. SCOPE predicts distributions of future occupancy grids from robot motion, tracked objects, and a local static map and inserts their mean and uncertainty into the costmaps of existing planners, with distillation enabling embedded inference~\cite{P_R18}. However, its navigation comparisons are mainly simulated, while its constant-velocity robot model fails during abrupt rotation. Homotopy-driven MPC optimizes distinct trajectory classes in parallel from externally tracked poses~\cite{P_S09}. Candidate diversity reduces local inconsistency without compensating for a missed pedestrian or an incorrect interaction model.

Learned local policies encode observation history implicitly. Visual proximal policy optimization (PPO) adds a motion-perceptibility auxiliary task~\cite{P_M04}, depth-based recurrent PPO is compared with DWA and TEB in simulation~\cite{P_M06}, and a LiDAR actor--critic controller uses spatiotemporal attention under adaptive Monte Carlo localization (AMCL) on a prebuilt map~\cite{P_M22}. HEIGHT separates pedestrians from static obstacles in a heterogeneous spatio-temporal graph and transfers from a low-fidelity simulator to four everyday environments, although its deployed obstacle input assumes a map and reliable localization~\cite{P_R16}. A rough-terrain sim-to-real policy trained with terrain scaling and pose noise still fails under vibration and inaccurate reconstructed geometry, with offboard inference over Wi-Fi~\cite{P_M42}. Lifelong navigation, in turn, learns only where an initial DWA policy supplies useful experience~\cite{P_M45}. As the learned-navigation review notes~\cite{P_X04}, the training data and observation representation determine which disturbances such controllers can accommodate. Recurrence alone therefore does not establish robustness to localization failure.

\subsection{Safety Filters and Explicit Risk Conditions}

A safety filter modifies a nominal command when a constraint would otherwise be violated. A LiDAR-based CBF controller fits a GP to current obstacle edges and uses its mean and gradient in a quadratic program, with an invariance argument that requires known state, unconstrained inputs, and suitable updates while excluding GP variance~\cite{P_E26}. A circulation-embedded CBF adds a tangential constraint and hold logic~\cite{P_E22}, dynamic movement primitives receive a CBF correction to their forcing term~\cite{P_E23}, an adaptive high-order CBF filters iterative LQR commands under input bounds~\cite{P_E24}, a disturbance-estimation CBF compensates matched uncertainty on a tracked robot~\cite{P_E29}, and safe reinforcement learning learns state-dependent high-order barrier penalties alongside soft actor--critic~\cite{P_E30}. Their evidence ranges from external pose and offboard computation to numerical studies. High-order tuning also affects both reaction time and feasibility under actuator bounds.

Other controllers constrain an explicit risk event. A distributionally robust CBF uses noisy barrier samples~\cite{P_M03}, a conditional value at risk (CVaR) controller selects among feasible risk levels in crowds~\cite{P_E27}, and a risk-aware-zone barrier expands an avoidance region under input saturation with a score that is not a calibrated probability~\cite{P_S05}. Safe Horizon MPC samples obstacle trajectories, applies a plan only after slack and support checks, and otherwise decelerates~\cite{P_E25}. In 100 Gaussian eight-pedestrian simulations with a joint-risk specification of 0.05, its mean traversal takes 16.0\,s (standard deviation 1.9\,s) versus 18.4\,s (2.4\,s) for chance-constrained MPC, with maximum estimated plan collision probabilities of 0.0092 and 0.0019. Faster traversal thus accompanies a less conservative plan under the assumed distribution, although the reported probabilities are not empirical collision rates. Perceive With Confidence calibrates obstacle inflation over sampled environments for a planner with an inevitable-collision-set constraint on a legged robot~\cite{P_F08}, whereas conformal crowd navigation limits cumulative intrusion into calibrated pedestrian regions on a Mecanum robot with offboard computation~\cite{P_S48}. MonoMPC learns the distribution of minimum clearance for candidate controls from monocular depth and minimizes a distribution-based risk~\cite{P_R14}. Jackal trials report up to seven times fewer collisions than a navigation stack on the same estimated depth, with planning on a networked laptop and no temporal memory. Chance constraints, CVaR, conformal coverage, and calibrated clearance concern different events. Comparisons should preserve these distinctions rather than read every confidence value as a collision probability.

\subsection{Tracking, Mission Costs, and Network Support}

Tracking control shows how model error, command age, and mission utility enter execution. A GP-MPC controller learns residual dynamics offline from field data and uses only the GP mean~\cite{P_M33}, whereas predictive learning refines repetitive Mecanum trajectories under external UWB pose~\cite{P_M09}. Delay-aware MPC predicts current states from delayed packets and buffers future controls~\cite{P_M27}, whereas vehicle--cloud MPC runs onboard and remote optimizers in parallel and stops on timeout~\cite{P_E21}, two different responses to communication loss. Mission-dependent costs add another reason to change a route, as in infrastructure-based hazard inference added to a cost map~\cite{P_E28} and in payload- or sensing-aware MPPI on a compressed scouting map~\cite{P_X12}. Prescribed-time spatiotemporal tubes assume bounded disturbances and terminal separation without explicit input constraints~\cite{P_N12}, memory-event-triggered MPC reduces repeated solves for a no-slip Ackermann model in hardware-in-the-loop tests~\cite{P_N14}, curvature-constrained vector fields co-design motion and saturation on an Ackermann vehicle with motion capture~\cite{P_S07}, and universal constrained tracking filters references that exceed prescribed state bounds~\cite{P_S08}. Neither a reach-avoid proof without input limits nor fewer solves establishes timely, feasible control on a saturated onboard system.

\subsection{Synthesis}

Motion methods differ in what they change, namely the terrain cost, the sampled route, the predicted occupancy or obstacle motion, or the final command. A terrain cost can favor better support without enforcing stability, whereas a learned policy can avoid familiar interactions without identifying a feasible response to an unseen fault. Explicit constraints make admissibility testable only under their state, dynamics, and update assumptions. The resulting trade-offs concern progress as well as conservatism, and comparisons should therefore establish which of these changes produces the reported task benefit.
